% !PRINTED 188
\weberpage{188}

Such figures are the regular pyramids, the double pyramids
and the regular solids (tetrahedron, octahedron, cube,
dodecahedron and icosahedron).

The regular pyramid of $g$ sides is brought into coincidence with
itself by rotation about the principal axis through an angle
$2\pi:g$ and, by repetition of this rotation, in $g$ different ways.

If $g$ is even, a regular double pyramid of $g$ side faces can be
brought into coincidence with itself not only by rotation about its
principal axis through the angle $4\pi:g$, but also (in $\tfrac{1}{2}g$
different ways) by rotation about an axis lying in the equatorial
plane, with an angle of rotation of 180 degrees.

The tetrahedron admits 12 different positions in which it
occupies the same space.

To obtain them, denote the place of a vertex in the unit position
by 1. One then obtains three positions of the tetrahedron in which
the vertex 1 remains fixed. But one can bring any vertex to the
place of 1, whereby the number is quadrupled.

The same number is also found if one holds fast the plane of a
side face, in which case one again finds three positions of the
tetrahedron, and brings each side face into the initial plane.

Finally one can also count in such a way that one brings each
edge into coincidence with a fixed segment in two ways. One proceeds
in the same manner with the remaining regular solids. One thus finds
the degree of the group

\begin{quote}
equal to the product of the number of vertices by the number of the
edges or side faces meeting in a vertex, or

equal to the product of the number of side faces by the number of
sides or vertices of a bounding face, or

equal to twice the number of edges.
\end{quote}

Each of these counts yields

\begin{quote}
for the tetrahedron 12 positions,\\
for the octahedron and the cube 24 positions,\\
for the dodecahedron and icosahedron 60 positions.
\end{quote}
